%!TEX encoding = UTF8
%!TEX root = 0-notes.tex

\chapter{Séries numériques}
\label{chap:1}

\section{Préliminaires calculatoires}

% inégalité triangulaire
% manipulations d'inéquations
% etc...
% à compiler au fur et à mesure pour présenter au début en exos calculatoires.

\section{Sommes partielles}

\subsection{Somme téléscopique}

\subsection{Somme géométrique}

% cas particulier de téléscopique

\subsection{Somme des $n^k$}

% cas particulier de téléscopique

\section{Série numérique : étude de convergence}

% on peut ignorer un nombre fini de termes lors de l'étude de convergence

\subsection{Comparaison géométrique}

\subsection{Critères de Cauchy et de d'Alembert}

\subsection{Critère de condensation}

% motivation par la série harmonique

\subsection{Comparaison intégrale}

% série de Riemann 1/n^s
% série de Bertrand 1/n / log(n)


% motivations pour chap suites : 
% 
% comparaison o() O() et ~ :
% pour l'étude de convergence
% 
% Cauchy : 
% étude de la somme résiduelle pour montrer la convergence
% 
% Fonction et DL :
% étudier précisément sum <-> int
% et voir la somme téléscopique comme int d/dx pousse à étudier f(n) = u_n
% et son intégrale. Par exemple sum 1/n = sum ln(n+1) - ln(n) + ?
% or ln(n+1)-ln(n) = ln(1+1/n) = ln(1) + 1/n + 1/n² + ... + reste O(1/n^r).
% donc sum 1/n = ln(n) + gamma + c/n + ... et sum 1/n - ln(n) tend vers gamma, constante d'Euler.


\section*{Exercices supplémentaires}
\addcontentsline{toc}{section}{Exercices supplémentaires}



\exe{,difficulty=0}{
	Vrai ou faux ? Démontrer que chaque proposition est vraie ou donner un contre-exemple.
	\begin{enumerate}
		\item
		Si $u_n$ tend vers 0, alors $\sum u_n$ converge.
	\end{enumerate}
}{exe:VF-series}{
	\begin{enumerate}
		\item
		Faux : la série harmonique $u_n = 1/n$ donne un contre-exemple.
	\end{enumerate}
}

\exe{,difficulty=0}{
	Montrer le critère de Leibniz\footnote{Gottfried Wilhelm Leibniz (1646 -- 1716), polymathe allemand.} suivant.
	
	Soit $(u_n)$ une suite vérifiant :
	\begin{enumerate}
		\item
		$u_n$ est alternée, donc $u_n u_{n+1} \leq 0$ pour tout $n\in\N$ ;
		\item
		la série absolue $|u_n|$ est décroissante et tend vers 0.
	\end{enumerate}
	Alors la série $\sum u_n$ est convergente.
}{exe:critere-leibniz}{
	todo
}

%\section*{Solutions}
%\addcontentsline{toc}{section}{Solutions}
%
%\shipoutAnswer



